\section{Waterfall mechanics}\label{sec:waterfall}

The tranche cash flows follow the notional-depleting waterfall of
\autocite{AndersenSideniusBasu2003}. Let $\tau \in \{\text{equity},
\text{mezz}, \text{senior}\}$ index the tranches, with attach / detach
points $(a_\tau, d_\tau)$ and initial notionals $N_\tau^0 = P_0 (d_\tau -
a_\tau)$. Two waterfalls run in parallel.

\paragraph{Loss allocation.}
At every period $t_k$, the aggregate loss $L(t_k)$ defined in
\cref{eq:agg-loss} is split across the stack by the stop-loss payoff
\cref{eq:stop-loss}. The surviving tranche notional is $N_\tau(t_k) =
N_\tau^0 - P_0 \cdot L_\tau(t_k)$. \Cref{fig:waterfall} illustrates the
order in which a unit of aggregate loss flows across the stack: the equity
tranche is wiped out for $L = 0.25$, the mezzanine for $L = 0.60$, and the
senior absorbs only the residual beyond $0.60$.

\paragraph{Interest waterfall.}
Each period the net rent inflow $R_{t_k} \Delta t$ is distributed
sequentially. Senior receives its promised coupon
$c_\text{senior} \cdot N_\text{senior}(t_k) \cdot \Delta t$ first, capped
by the available cash; mezzanine receives the same on its own notional out
of what remains; equity collects whatever is left:
\begin{align*}
C_\text{senior}(t_k) &= \min\bigl( R_{t_k}\Delta t,\ c_\text{senior}
N_\text{senior}(t_k) \Delta t \bigr), \\
C_\text{mezz}(t_k) &= \min\bigl( R_{t_k}\Delta t - C_\text{senior}(t_k),\
c_\text{mezz} N_\text{mezz}(t_k) \Delta t \bigr), \\
C_\text{equity}(t_k) &= R_{t_k}\Delta t - C_\text{senior}(t_k) - C_\text{mezz}(t_k).
\end{align*}
The equity cash flow can be negative when defaults force the SAS to use
reserves; in our simulation we floor net rent at zero before this
allocation.

\paragraph{Principal waterfall.}
At the terminal date $T$ the building is sold at $V_T = S_T \cdot (1 -
L(T))$ --- market price scaled by the surviving fraction of the portfolio.
The proceeds repay senior first (up to $N_\text{senior}(T)$), then
mezzanine (up to $N_\text{mezz}(T)$), and equity collects any residual.
Equity therefore captures the property's capital appreciation \emph{net of
the cumulative writedown}; this is the upside that justifies its first-loss
status.

\paragraph{ASB recursion in matrix form.}
Let $N_\tau^{(k)}$ denote the notional of tranche $\tau$ at the end of
period $k$, with the convention $N_\tau^{(0)} = N_\tau^0$. The
notional-update rule is
\begin{equation*}
N_\tau^{(k)} = N_\tau^0 - P_0 \cdot L_\tau\bigl(L^{(k)}\bigr),
\quad L^{(k)} = \sum_{i: \tau_i \le t_k} \tfrac{\ell_i}{n},
\end{equation*}
where the per-tranche loss $L_\tau$ follows the stop-loss payoff
\cref{eq:stop-loss}. The interest waterfall at period $k$ is the
sequential allocation
\begin{equation*}
C_\tau^{(k)} = \min\bigl(c_\tau N_\tau^{(k)} \Delta t,\ R^{(k)} - {\textstyle\sum_{\tau' > \tau}} C_{\tau'}^{(k)}\bigr),
\quad \tau \in \{\text{senior}, \text{mezz}, \text{equity}\},
\end{equation*}
processed senior-to-junior. The principal waterfall at $T$ is the same
recursion applied to the surviving sale proceeds
$V_T = S_T(1 - L^{(N)})$, with the senior receiving back its remaining
notional first, then the mezzanine, then the equity collecting whatever
residual remains (which can include capital appreciation
$V_T > \sum_\tau N_\tau^0$ when $S_T$ has grown).

\paragraph{Fair coupon.}
The fair coupon $c_\tau^{\star}$ of a tranche is the value that sets its
present value equal to par at $t = 0$. We solve it for one tranche at a
time using Brent's method on the pricing residual; the other tranches keep
the contract coupons specified in \texttt{config/}. The fair-spread
solution is available in \texttt{pricing/tranche\_pricer.py::solve\_fair\_coupon}
and used in the supplementary sensitivity runs.

\begin{figure}[ht]
  \centering
  \includegraphics[width=0.85\linewidth]{fig_waterfall_explainer.pdf}
  \caption{Tranche-by-tranche loss absorption as the aggregate cumulative
  loss grows from zero to par. The dashed verticals mark the detach points
  beyond which each tranche is exhausted.}
  \label{fig:waterfall}
\end{figure}
